\documentclass[10pt,a4paper]{amsart}
\usepackage[margin=19mm]{geometry}
\usepackage{amsmath,amssymb,mathtools}
\usepackage[scr=rsfs,cal=euler]{mathalfa}
\usepackage{xfrac}
\usepackage[unicode,hidelinks,bookmarks=false]{hyperref}
\newcommand{\ie}{i.e.\ }
\newcommand{\cf}{cf.\ }
\newcommand{\resp}{respectively\ }
\newcommand{\Subsection}{\subsection}
\providecommand{\nobreakdash}{\nobreak\hbox{-}}
\setlength{\emergencystretch}{2em}
\multlinegap0pt
\usepackage{bm}
\usepackage{bbm}
\usepackage{dsfont}
\usepackage{xspace}
\usepackage{cancel}
\usepackage{mathtools}
\usepackage{cleveref}
\usepackage{amsthm}
\makeatletter
\@ifundefined{theorem}{\newtheorem{theorem}{Theorem}}{}
\@ifundefined{lemma}{\newtheorem{lemma}[theorem]{Lemma}}{}
\@ifundefined{prop}{\newtheorem{prop}[theorem]{Proposition}}{}
\makeatother
\DeclareMathOperator*{\argmin}{arg\,min}
\newcommand{\Lip}{\operatorname{Lip}}
\newcommand{\N}{\mathcal{N}}
\newcommand{\R}{\mathbb{R}}
\newcommand{\conv}{\operatorname{conv}}
\newcommand{\weakk}{\stackrel{\ast}{\rightharpoonup}}
\title[Learning Firmly Nonexpansive Operators]{Learning Firmly Nonexpansive Operators}
\author{Kristian Bredies, Jonathan Chirinos-Rodriguez, Emanuele Naldi}
\date{Source: arXiv:2407.14156v3 (CC BY 4.0)}
\begin{document}
\maketitle

\noindent\emph{Source and licence.} Learning Firmly Nonexpansive Operators. Version arXiv:2407.14156v3 (CC BY 4.0), distributed under the Creative Commons Attribution 4.0 International licence, as recorded in the arXiv licence metadata for this exact version and shown on the abstract page https://arxiv.org/abs/2407.14156v3. Full rights notice: licenses/arxiv-2407.14156-CC-BY-4.0.txt.\par\smallskip

\noindent\textit{Editorial scope.} Counted unit: the Theorem whose source label is \texttt{thm:density} in the article (source0.tex lines 851--853, source label \texttt{thm:density}), delivered together with the article's own complete original proof of it (source0.tex lines 854--931, one \texttt{proof} environment of 78 lines). No mathematics is written, repaired or re-derived: statement and proof are the article's own text line for line. Required context retained uncounted: eq:DP (lines 662--665); eq:matrices (lines 739--741); prop:aediff (lines 759--761); lem:bound\_invA (lines 807--811); lem:311 (lines 839--845); eq:Lk*bounded (lines 841--843). Every definition, display and label of those lines is retained unchanged. Omitted from this package and not redistributed: 1--661, 666--738, 742--758, 762--806, 812--838, 844--850, 932--1323. Nothing inside a retained excerpt is omitted or altered: every retained body is delivered exactly as the article's lines are written, so no omission receipt of any kind accompanies this package. Citations: the retained excerpts carry no citation command at all, so no bibliography entry of the article is transcribed and no reference is redistributed. Reference adaptation: none was needed and none was made; every reference inside the delivered unit resolves inside it, so no unresolved reference or citation is produced. Printed slips kept literal: no printed word, symbol, hypothesis or number of the counted statement or of its proof was altered. Layout and class: the article's own class is \texttt{siamonline250211} with packages including lipsum, graphicx, stackengine, amsfonts, amssymb, epstopdf, algorithmic, comment, dsfont, hyperref, enumitem, amsopn; the wrapper is the lane's pinned \texttt{amsart} editorial wrapper at 10pt on A4, which restates the theorem-like environments the retained text needs and adds the article's own macro definitions that the retained text needs (\texttt{\textbackslash Lip}, \texttt{\textbackslash N}, \texttt{\textbackslash R}, \texttt{\textbackslash argmin}, \texttt{\textbackslash conv}, \texttt{\textbackslash weakk}), with extra wrapper packages (bm, bbm, dsfont, xspace, cancel, mathtools, cleveref). The delivered numbering is the unit's own. Assets: no retained span contains a figure environment, a graphics-inclusion command, a PDF-page inclusion, a table or a TikZ picture; no image file is copied into this package, the package declares no asset dependency and no image file is redistributed with it. Licence: version arXiv:2407.14156v3 of this article is distributed under the Creative Commons Attribution 4.0 International licence, as recorded in the arXiv licence metadata for this exact version and shown on the abstract page https://arxiv.org/abs/2407.14156v3. A full rights notice is delivered at licenses/arxiv-2407.14156-CC-BY-4.0.txt. Characters: every retained excerpt is pure ASCII, so the delivered engine, XeLaTeX with its default Latin Modern text font, reports no missing character.
\begin{equation}\label{eq:DP}
\tag{DP}
\widehat{N}\in \argmin_{N\in\N}\widehat{F}(N), \quad \widehat{F}(N):=\frac{1}{n}\sum_{i=1}^n \|N(\bar{x}_i)-\bar{y}_i\|^2.
\end{equation}

\begin{equation}\label{eq:matrices}
A_t:=[x_{i_1}-x_{i_0}\mid\dots \mid x_{i_d}-x_{i_0}], \quad B_t=[y_{i_1}-y_{i_0}\mid\ldots \mid y_{i_d}-y_{i_0}].
\end{equation}

\begin{prop}\label{prop:aediff}
Let $N$ be a nonexpansive operator. Then, it is differentiable almost everywhere and, for every differentiability point $x$, it holds that $\|\nabla N(x)\|\leq 1$.
\end{prop}

\begin{lemma}\label{lem:bound_invA} Let $\mathfrak{F}=(\mathfrak{T}_k)_{k\in\mathbb{N}}$ be a strongly regular vanishing family of simplicial partitions for $\mathrm{conv}(D)$. Then, there exists a constant $c'>0$ such that, for every partition $\mathfrak{T}_k$ in $\mathfrak{F}$ and all simplices $S_t\in \mathfrak{T}_k$, $t=1,\ldots, \, \ell(k)$, $\ell(k)\in\mathbb{N}$, if we denote by $A_t$ the matrix  relative to the simplex $S_t$ defined in \cref{eq:matrices}, we have
$$
\|A_t^{-1}\|\leq \frac{k}{c'}.
$$
\end{lemma}

\begin{lemma}\label{lem:311}
Let $\mathfrak{F}=(\mathfrak{T}_k)_{k\in\mathbb{N}}$ be a vanishing family of simplicial partitions for $\conv(D)$ and define, for every $k\in\mathbb{N}$,
\begin{equation}\label{eq:Lk*bounded}
\widehat{L}_k\in\argmin_{N\in \mathrm{PA}(\mathfrak{T}_k)\cap\mathcal{N}}\widehat{F}(N),
\end{equation}
where $\widehat{F}$ is defined as in \cref{eq:DP}. Then, the sequence $(\widehat{L}_k)_{k\in\mathbb{N}}$ is bounded in $\Lip(\R^d, \R^d)$ and there exists a subsequence $(\widehat{L}_{k_j})_{j\in\mathbb{N}}$ and an operator $\widehat{L}\in \mathcal{N}$ such that $\widehat{L}_{k_j}\weakk \widehat{L}$ as $j \to \infty$.
\end{lemma}

\begin{theorem}\label{thm:density} Let $\mathfrak{F}=(\mathfrak{T}_k)_{k\in \mathbb{N}}$ be a strongly regular vanishing family of simplicial partitions for $\mathrm{conv}(D)$. Then, for every sequence of operators defined as in \cref{eq:Lk*bounded}, there exists a subsequence $(\widehat{L}_{k_j})_{j\in\mathbb{N}}$ such that
\[ \widehat{L}_{k_j}\weakk \widehat{L} \ \text{as} \ j \to \infty \ \ \text{ with } \ \ \widehat{L}\in\argmin_{N\in\mathcal{N}}\widehat{F}(N).\]
\end{theorem}

\begin{proof}
For every $k\in\mathbb{N}$, let $(x_i^k)_{i=1}^{n_k}$ be the vertices of all the simplices of the partition $\mathfrak{T}_k$ and recall that $\bar x_i=x_i^k$ for each $i=1,\ldots, n$ and every $k\in\mathbb{N}$. Let $\widehat{N}$ be a minimizer of \cref{eq:DP}, $\delta > 0$ and define the operator $N^\delta:=\widehat{N}\ast G_\delta$, where $G_\delta:\R^d\to\R$ is a mollifier; i.e., defined for every $x\in\R^d$ as
$$
G_\delta(x):=\frac{1}{\delta^d}G\left(\frac{x}{\delta}\right), \quad G\in \mathcal{C}^\infty(\R^d),\quad \mathrm{supp} \, G\subset B_1(0),\quad G\geq 0, \quad \int_{\R^d} G(x) \, \mathrm{d}x=1.
$$
By \cref{prop:aediff}, $\widehat{N}$ is almost everywhere differentiable and thus, $\|\nabla N^\delta\|_\infty\leq 1$. Indeed, by $G_\delta \geq 0$ and $\| G_\delta \|_1 = 1$,
$$
\begin{aligned}
\|\nabla N^\delta\|_\infty &=\sup_{x\in\R^d} \left\|\int_{\R^d}\nabla \widehat{N}(x-y) G_\delta(y) \, \mathrm{d}y \right\| \\
&\leq \sup_{x\in\R^d}\int_{\R^d} \left\| \nabla \widehat{N}(x-y) \right\| G_\delta(y) \, \mathrm{d}y\\
&\leq \|\nabla \widehat{N}\|_\infty\|G_\delta\|_1\leq 1.
\end{aligned}
$$
Similarly, we have $\|N^\delta - \widehat{N}\|_\infty \leq \delta$ since for each $x \in \mathbb{R}^d$, we can estimate 
\[
\|N^\delta(x) - \widehat{N}(x)\| \leq \int_{\mathbb{R}^d} \|\widehat{N}(x - y) - \widehat{N}(x)\| G_\delta(y) \, \mathrm{d}y
\leq \delta \|G_\delta\|_1 \leq \delta
\]
as $\mathcal{L}(\widehat{N}) \leq 1$ and $G_\delta$ is supported on $B_\delta(0)$.

Additionally, since $\nabla N^\delta$ is uniformly continuous on $\conv (D)$, there exists $\varepsilon>0$ such that
\begin{equation}\label{eq:312unifcont}
\|N^\delta(x)-N^\delta(x')-\nabla N^\delta(x')(x-x')\|\leq \delta\|x-x'\|
\end{equation}
for every $x, x'\in\conv(D)$ with $\|x-x'\|\leq \varepsilon$. Next, for each $k\in\mathbb{N}$, let  $L_k^\delta$ be the unique mapping in $\mathrm{PA}(\mathfrak{T}_k)$ such that $L_k^\delta(x_i^k)=N^\delta(x^k_i)$ for $i=1,\ldots, n_k$. For each $S_t\in\mathfrak{T}_k$ we denote by $(x^k_{i(t,j)})_{j=0}^d$ its vertices, and recall that the matrices $A_t$ and $B_t$ are given by
$$
A_t=[x^k_{i(t,1)}-x^k_{i(t,0)}\mid\dots \mid x^k_{i(t,d)}-x^k_{i(t,0)}],
$$
and 
$$
B_t=[N^\delta(x^k_{i(t,1)})-N^\delta(x^k_{i(t,0)})\mid\ldots \mid N^\delta(x^k_{i(t,d)})-N^\delta(x^k_{i(t,0)})].
$$
By assumption, $\|x^k_{i(t,j)}-x^k_{i(t,0)}\|\leq 1/k$ for each $S_t\in\mathfrak{T}_k$ and $j=1,\dots d$. Therefore, by choosing $k\in\mathbb{N}$ large enough so that $(1/k)\leq\varepsilon$, \cref{eq:312unifcont} gives
\begin{equation*}
\|N^\delta(x^k_{i(t,j)})-N^\delta(x^k_{i(t,0)})-\nabla N^\delta(x^k_{i(t,0)})(x^k_{i(t,j)}-x^k_{i(t,0)})\|\leq \delta\|x^k_{i(t,j)}-x^k_{i(t,0)}\| \leq \frac{\delta}{k},
\end{equation*}
for $j=1,\ldots, d$, and so
\begin{multline*}
\|B_t-\nabla N^\delta(x^k_{i(t,0)})A_t\| \leq 
\|B_t-\nabla N^\delta(x^k_{i(t,0)})A_t\|_{\mathrm{F}} \\
= \Biggl( \sum_{j=1}^d \| N^\delta(x^k_{i(t,j)})-N^\delta(x^k_{i(t,0)})-\nabla N^\delta(x^k_{i(t,0)})(x^k_{i(t,j)}-x^k_{i(t,0)})  \|^2 \Biggr)^{1/2} \leq
\frac{\sqrt{d}\delta}{k}.    
\end{multline*}
Consequently, by \cref{lem:bound_invA}, we obtain
$$
\begin{aligned}
\|B_tA_t^{-1}\|&\leq \|\nabla N^\delta(x^k_{i(t,0)})\|+\|B_tA_t^{-1}- \nabla N^\delta(x^k_{i(t,0)})\|\\
&\leq 1+ \|B_t - \nabla N^\delta(x_{i(t,0)}) A_t \| \| A_t^{-1} \| \leq 1+c\delta,
\end{aligned}
$$
for a suitable constant $c>0$. Since $\|B_tA_t^{-1}\|$ is
the norm of the gradient of $L_k^\delta$ in each simplex $S_t$, setting
$$
\widetilde{L}^\delta_k=\frac{1}{1+c\delta}L^\delta_k,
$$
yields $\widetilde{L}^\delta_k\in\mathrm{PA}(\mathfrak{T}_k)\cap \N$; i.e., $\widetilde{L}^\delta_k$ is feasible for \cref{eq:Lk*bounded}. By recalling that, for every $i=1,\ldots, n$, it holds $L^\delta_k(\bar x_i)=N^\delta(\bar x_i)$ and $\|N^\delta-\widehat{N}\|_\infty\leq \delta$, we have
$$
\begin{aligned}
\|\widetilde{L}_k^\delta(\bar x_i)-\widehat{N}(\bar x_i)\|&\leq \|\widetilde{L}_k^\delta(\bar x_i)- L_k^\delta(\bar x_i)\|+\|N^\delta(\bar x_i)-\widehat{N}(\bar x_i)\|\\
&\leq c\delta \bigl( \| \widehat{N}(\bar x_i) \| + \|N^\delta(\bar x_i) - \widehat{N}(\bar x_i) \| \bigr) + \delta\\
&\leq (1+c M+c\delta)\delta, 
\end{aligned}
$$
for $i=1,\ldots, n$, where $M:= \sup_{x \in \conv(D)} \| \widehat{N}(x) \| < \infty$. If $\widehat{L}_k$ minimizes \cref{eq:Lk*bounded} for each $k\in\mathbb{N}$, it follows that
\begin{align*}
\widehat{F}(\widehat{L}_k)&\leq \widehat{F}(\widetilde{L}_k^\delta)\leq \frac1n\sum_{i=1}^n(1+\delta)\|\widehat{N}(\bar x_i)-\bar y_i\|^2+\left(1+\frac1\delta\right)\|\widetilde{L}_k^\delta(\bar x_i)-\widehat{N}(\bar x_i)\|^2\\
&\leq (1+\delta)\widehat{F}(\widehat{N})+\left(\delta^2+\delta\right)(1+c M+c\delta)^2,
\end{align*}
where we have used Young's inequality and the definition of $\widehat{F}$. As $\delta>0$ was arbitrary, we can conclude that
$$
\limsup_{k\to\infty} \widehat{F}(\widehat{L}_k)\leq \widehat{F}(\widehat{N}).
$$
By \cref{lem:311}, there exists a subsequence  $(\widehat{L}_{k_j})_{j\in\mathbb{N}}$ of $(\widehat{L}_k)_{k\in\mathbb{N}}$ such that $\widehat{L}_{k_j}\weakk \widehat{L}$ as $j\to\infty$ for some $\widehat{L} \in \N$. The weak$^*$ lower semi-continuity of $\widehat{F}$ then gives
$$
\widehat{F}(\widehat{L})\leq \liminf_{j\to\infty}\widehat{F}(\widehat{L}_{k_j})\leq \limsup_{k\to\infty}\widehat{F}(\widehat{L}_k)\leq \widehat{F}(\widehat{N}).
$$
As we have shown that $\widehat{L}$ is a minimizer of \cref{eq:DP}, we conclude.
\end{proof}

\end{document}
